% !PRINTED 107
\weberpage{107}

\begin{center}
{\large \S.\ 32.}

\vspace{0.5em}
{\large\gesperrt{Vorzeichenwechsel von $f(x)$. Wurzeln von Gleichungen ungeraden Grades und von reinen Gleichungen.}}
\end{center}

Auf Grund der Sätze des vorigen Paragraphen können wir das folgende wichtige Theorem beweisen.

\medskip
\noindent VIII. \gesperrt{Sind die Coëfficienten von $f(x)$ reell und existiren zwei reelle Werthe $a,b$ von $x$, so dass $f(a)$ und $f(b)$ entgegengesetzte Vorzeichen haben, so giebt es zwischen $a$ und $b$ wenigstens eine Wurzel $\xi$ der Gleichung $f(x)=0$.}
\medskip

Wir wollen annehmen, es sei $a<b$ und $f(a)$ negativ, $f(b)$ positiv. Wir theilen die Zahlen zwischen $a$ und $b$ so in zwei Theile $\mathfrak A$ und $\mathfrak B$, dass eine Zahl $\alpha$ zu $\mathfrak A$ gehört, wenn zwischen $\alpha$ und $a$, die Grenzen eingerechnet, kein $x$ liegt, wofür $f(x)$ positiv wird, und eine Zahl $\beta$ zu $\mathfrak B$, wenn zwischen $\beta$ und $a$ wenigstens ein Werth von $x$ liegt, für den $f(x)$ positiv wird. Ein drittes ist offenbar nicht möglich und jede Zahl zwischen $a$ und $b$ ist also in $\mathfrak A$ oder in $\mathfrak B$ untergebracht. Wir rechnen noch die Zahlen, die kleiner als $a$ sind zu $\mathfrak A$, und die grösser als $b$ sind zu $\mathfrak B$.

Gehört $\alpha$ zu $\mathfrak A$, so gehört auch jedes $x$, was kleiner ist als $\alpha$ zu $\mathfrak A$, und gehört $\beta$ zu $\mathfrak B$, so gehört jedes grössere $x$ zu $\mathfrak B$. Es ist also jedes $\beta$ grösser als jedes $\alpha$ und die beiden Zahlengebiete $\mathfrak A$, $\mathfrak B$ bilden einen \gesperrt{Schnitt} (Einleitung, S. 5) der durch eine Zahl $\xi$ erzeugt wird, so dass jede Zahl $\alpha$, die kleiner als $\xi$ ist, zu $\mathfrak A$ gehört, jede Zahl $\beta$, die grösser als $\xi$ ist, zu $\mathfrak B$.

Ist nun für irgend eine Zahl $x_0$ der Werth von $f(x_0)$ negativ oder positiv, so lässt sich nach V. des vorigen Paragraphen eine positive Grösse $\varepsilon$ so bestimmen, dass für jedes $x$ zwischen $x_0-\varepsilon$ und $x_0+\varepsilon$ die Werthe von $f(x)$ immer negativ oder immer positiv sind. Daraus ergiebt sich, dass $f(\xi)$ weder negativ noch positiv sein kann und also Null sein muss. Denn wäre $f(\xi)$ negativ, so wäre auch $f(x)$ zwischen $\xi$ und $\xi+\varepsilon$ negativ; diese Werthe würden also zu $\mathfrak A$ gehören, während sie doch grösser als $\xi$ sind; und wäre $f(\xi)$ positiv, so wäre $f(x)$ positiv, so lange $x$ zwischen $\xi-\varepsilon$ und $\xi$ liegt; diese Werthe würden also, obwohl sie kleiner als $\xi$ sind, zu $\mathfrak B$ gehören und beides ist nach der Definition von
%CONT>>
